\cotefile{84}

\scanpage{35}

\uline{Dém} OPS $A=O[u]/(u^2+bu+c)$, et un \unsure{endom.|autom.} $\Theta$ de la forme envisagée doit transformer $u$ en \deleted{\unsure{et}}

\[
u'=\bar u+\alpha.1_A \qquad \alpha\in\Gamma\underline O
\]

(puisque $u'\equiv\bar u\equiv-u\bmod\underline O.1_A$), avec

\[
\left\{\begin{aligned}
\operatorname{Tr}(u')&=\operatorname{Tr}(u) (=-b)\\
N(u')&=N(u)(=c)
\end{aligned}\right.
\]

--- les $\Theta$ en question correspondent biunivoquement aux $u'$ (sans condition sur $A$)$^*$.

% slanted
\margin{Sauf la condition pour $\Theta$ en outre \unsure{pour|que} \unsure{l'être|n'être|être} l'identité sur un \deleted{\unsure{ouv.}} \unsure{un vois.} de $X$, qu'il faut regarder à part, en un point $s\in X_0\ldots$}

En termes de $\alpha$, les conditions s'écrivent respectivement

\[
\left\{\begin{aligned}
2\alpha&=0 &&(\text{car }\operatorname{Tr}\alpha.1_A=2\alpha)\\
\alpha(\alpha-b)&=0
\end{aligned}\right.
\]

En les pts de $X-X_0$, \deleted{\unsure{le}} \inserted{$2\alpha=0$} implique $\alpha=0$.

Reste à regarder en un point $s$ de $X_0$. \deleted{\unsure{Je} \unsure{dis}} \inserted{je dis que $\alpha_s\in M_s=\operatorname{rad}\underline O_{X,s}$ --- sinon,} \deleted{par le corps résiduel $k(s)$, on aura}

\[
\deleted{\unsure{\text{que}|\text{donc}}\qquad\alpha(s)(\alpha(s)-b(s))=0,}
\]

\deleted{mais par hypothèse \deleted{$b(s)$} $A(s)$ étale i.e. $b(s)\ne 0$, donc on doit avoir ou bien $\alpha(s)=0$, ou bien $\alpha(s)=b(s)$,} $\alpha_s$ serait inversible, et on aurait, à cause de \inserted{la relation} $\alpha(\alpha-b)=0$, $\alpha_s-b_s=0$ i.e. $\alpha_s=b_s$, \deleted{\unsure{Ma|Mais}} i.e. $\alpha=b$ au voisinage de $s$, \deleted{donc} et comme $2\alpha=0$, on aurait $2b=0$, donc comme $b_s$ est \uline{inversible} (car $A$ étale en $s$, et $s\in X_0$) $2.\deleted{\unsure{\text{sera}|\text{serait}}}1_{X}$ serait nul au vois. de $s$ i.e. $X$ de car. $2$.

% archive sheet number: 34
